\documentclass[12pt,letterpaper]{article}
\usepackage[letterpaper,margin=0in]{geometry}
\usepackage[T1]{fontenc}
\usepackage{lmodern}
\usepackage{tikz}
\usetikzlibrary{calc}
\pdfmapfile{+lm.map}
\pagestyle{empty}
\hyphenpenalty=10000
\exhyphenpenalty=10000
\setlength{\parindent}{0pt}
\begin{document}
\noindent\begin{tikzpicture}[remember picture,overlay,x=1in,y=1in,shift={(current page.south west)}]
\node[anchor=north west,inner sep=0pt,align=left,text width=7.18in,font=\fontsize{13}{17}\selectfont] at (0.64,10.48) {Week 20 / Averaging / Grades 4--5};
\draw[line width=.5pt] (.64,10.13)--(7.86,10.13);
\node[anchor=north west,inner sep=0pt,align=left,text width=7.18in,font=\fontsize{9.5}{13.5}\selectfont] at (0.64,0.48) {Bellingham Math Circle / Week 20 / GA20-45-v3};
\node[anchor=north east,inner sep=0pt,font=\fontsize{9.5}{12}\selectfont] at (7.86,.48) {1};
\node[anchor=north west,inner sep=0pt,align=left,text width=7.18in,font=\fontsize{12.5}{16.5}\selectfont] at (0.64,9.86) {Add the numbers joined to a circle and divide by how many numbers you added. The answer must equal the circle's number. Squares keep their numbers. Every circle must work at the same time.};
\node[anchor=north west,inner sep=0pt,align=left,text width=7.18in,font=\fontsize{14.5}{18.5}\selectfont] at (0.64,6.87) {\textbf{Problem 1:} Fill all the circles. Try to find a second filling for each board while keeping its square numbers.};
\node[anchor=north west,inner sep=0pt,align=left,text width=7.18in,font=\fontsize{12.5}{16.5}\selectfont] at (0.64,8.98) {Example: both circles work.};
\node[draw,fill=white,rectangle,minimum size=.60in,line width=.8pt,inner sep=0pt,font=\fontsize{16}{20}\selectfont] (ex0) at (1.4,8.32) {2};
\node[draw,fill=white,circle,minimum size=.60in,line width=.8pt,inner sep=0pt,font=\fontsize{16}{20}\selectfont] (ex1) at (3.3,8.32) {3};
\node[draw,fill=white,circle,minimum size=.60in,line width=.8pt,inner sep=0pt,font=\fontsize{16}{20}\selectfont] (ex2) at (5.199999999999999,8.32) {4};
\node[draw,fill=white,rectangle,minimum size=.60in,line width=.8pt,inner sep=0pt,font=\fontsize{16}{20}\selectfont] (ex3) at (7.1,8.32) {5};
\draw[line width=.85pt] (ex0)--(ex1);
\draw[line width=.85pt] (ex1)--(ex2);
\draw[line width=.85pt] (ex2)--(ex3);
\node[inner sep=.03in,font=\fontsize{13}{17}\selectfont] (three) at (3.3,7.62) {$(2+4)/2=3$};
\draw[->,line width=.55pt] (three.north)--(ex1.south);
\node[inner sep=.03in,font=\fontsize{13}{17}\selectfont] (four) at (5.2,7.62) {$(3+5)/2=4$};
\draw[->,line width=.55pt] (four.north)--(ex2.south);
\node[anchor=north west,inner sep=0pt,align=left,text width=7.18in,font=\fontsize{11.5}{15.5}\selectfont] at (0.64,7.2700000000000005) {Check both circles on this same unchanged board.};
\node[draw,fill=white,rectangle,line width=1.05pt,minimum size=0.68in,inner sep=0pt,font=\fontsize{19}{25}\selectfont] (g10) at (1.4,5.9) {0};
\node[draw,fill=white,circle,line width=1.05pt,minimum size=0.68in,inner sep=0pt,font=\fontsize{19}{25}\selectfont] (g11) at (3.3,5.9) {};
\node[draw,fill=white,circle,line width=1.05pt,minimum size=0.68in,inner sep=0pt,font=\fontsize{19}{25}\selectfont] (g12) at (5.2,5.9) {};
\node[draw,fill=white,rectangle,line width=1.05pt,minimum size=0.68in,inner sep=0pt,font=\fontsize{19}{25}\selectfont] (g13) at (7.1,5.9) {15};
\draw[line width=1.15pt] (g10)--(g11);
\draw[line width=1.15pt] (g11)--(g12);
\draw[line width=1.15pt] (g12)--(g13);
\node[draw,fill=white,rectangle,line width=1.05pt,minimum size=0.68in,inner sep=0pt,font=\fontsize{19}{25}\selectfont] (g2a) at (2.05,5.05) {0};
\node[draw,fill=white,rectangle,line width=1.05pt,minimum size=0.68in,inner sep=0pt,font=\fontsize{19}{25}\selectfont] (g2c) at (2.05,3.45) {8};
\node[draw,fill=white,circle,line width=1.05pt,minimum size=0.68in,inner sep=0pt,font=\fontsize{19}{25}\selectfont] (g2u) at (3.4499999999999997,4.25) {};
\node[draw,fill=white,circle,line width=1.05pt,minimum size=0.68in,inner sep=0pt,font=\fontsize{19}{25}\selectfont] (g2v) at (4.75,4.25) {};
\node[draw,fill=white,rectangle,line width=1.05pt,minimum size=0.68in,inner sep=0pt,font=\fontsize{19}{25}\selectfont] (g2b) at (6.05,4.25) {9};
\draw[line width=1.15pt] (g2a)--(g2u);
\draw[line width=1.15pt] (g2c)--(g2u);
\draw[line width=1.15pt] (g2u)--(g2v);
\draw[line width=1.15pt] (g2v)--(g2b);
\node[draw,fill=white,rectangle,line width=1.05pt,minimum size=0.68in,inner sep=0pt,font=\fontsize{19}{25}\selectfont] (g3a) at (1.9500000000000002,1.9) {2};
\node[draw,fill=white,rectangle,line width=1.05pt,minimum size=0.68in,inner sep=0pt,font=\fontsize{19}{25}\selectfont] (g3b) at (6.55,1.9) {14};
\node[draw,fill=white,circle,line width=1.05pt,minimum size=0.68in,inner sep=0pt,font=\fontsize{19}{25}\selectfont] (g3u) at (4.25,2.6799999999999997) {};
\node[draw,fill=white,circle,line width=1.05pt,minimum size=0.68in,inner sep=0pt,font=\fontsize{19}{25}\selectfont] (g3v) at (4.25,1.1199999999999999) {};
\draw[line width=1.15pt] (g3a)--(g3u);
\draw[line width=1.15pt] (g3a)--(g3v);
\draw[line width=1.15pt] (g3b)--(g3u);
\draw[line width=1.15pt] (g3b)--(g3v);
\draw[line width=1.15pt] (g3u)--(g3v);
\end{tikzpicture}\newpage
\noindent\begin{tikzpicture}[remember picture,overlay,x=1in,y=1in,shift={(current page.south west)}]
\node[anchor=north west,inner sep=0pt,align=left,text width=7.18in,font=\fontsize{13}{17}\selectfont] at (0.64,10.48) {Week 20 / Averaging / Grades 4--5};
\draw[line width=.5pt] (.64,10.13)--(7.86,10.13);
\node[anchor=north west,inner sep=0pt,align=left,text width=7.18in,font=\fontsize{9.5}{13.5}\selectfont] at (0.64,0.48) {Bellingham Math Circle / Week 20 / GA20-45-v3};
\node[anchor=north east,inner sep=0pt,font=\fontsize{9.5}{12}\selectfont] at (7.86,.48) {2};
\node[anchor=north west,inner sep=0pt,align=left,text width=7.18in,font=\fontsize{14.5}{18.5}\selectfont] at (0.64,9.86) {\textbf{Problem 2:} Can either board have a circle number of 13 or more while all circles follow the rule? Find a filling, or explain why none can work.};
\node[draw,fill=white,rectangle,line width=1.05pt,minimum size=0.68in,inner sep=0pt,font=\fontsize{19}{25}\selectfont] (g4a) at (1.3,7.6) {0};
\node[draw,fill=white,circle,line width=1.05pt,minimum size=0.68in,inner sep=0pt,font=\fontsize{19}{25}\selectfont] (g4u) at (2.85,7.6) {};
\node[draw,fill=white,circle,line width=1.05pt,minimum size=0.68in,inner sep=0pt,font=\fontsize{19}{25}\selectfont] (g4v) at (4.4,7.6) {};
\node[draw,fill=white,rectangle,line width=1.05pt,minimum size=0.68in,inner sep=0pt,font=\fontsize{19}{25}\selectfont] (g4b) at (5.95,7.6) {6};
\node[draw,fill=white,circle,line width=1.05pt,minimum size=0.68in,inner sep=0pt,font=\fontsize{19}{25}\selectfont] (g4w) at (4.4,5.9) {};
\node[draw,fill=white,circle,line width=1.05pt,minimum size=0.68in,inner sep=0pt,font=\fontsize{19}{25}\selectfont] (g4z) at (6.25,5.9) {};
\node[draw,fill=white,rectangle,line width=1.05pt,minimum size=0.68in,inner sep=0pt,font=\fontsize{19}{25}\selectfont] (g4c) at (7.35,4.55) {11};
\draw[line width=1.15pt] (g4a)--(g4u);
\draw[line width=1.15pt] (g4u)--(g4v);
\draw[line width=1.15pt] (g4v)--(g4b);
\draw[line width=1.15pt] (g4v)--(g4w);
\draw[line width=1.15pt] (g4w)--(g4z);
\draw[line width=1.15pt] (g4z)--(g4v);
\draw[line width=1.15pt] (g4z)--(g4c);
\node[draw,fill=white,rectangle,line width=1.05pt,minimum size=0.68in,inner sep=0pt,font=\fontsize{19}{25}\selectfont] (g5a) at (1.9500000000000002,2.65) {2};
\node[draw,fill=white,rectangle,line width=1.05pt,minimum size=0.68in,inner sep=0pt,font=\fontsize{19}{25}\selectfont] (g5b) at (6.55,2.65) {10};
\node[draw,fill=white,circle,line width=1.05pt,minimum size=0.68in,inner sep=0pt,font=\fontsize{19}{25}\selectfont] (g5u) at (4.25,3.65) {};
\node[draw,fill=white,circle,line width=1.05pt,minimum size=0.68in,inner sep=0pt,font=\fontsize{19}{25}\selectfont] (g5v) at (4.25,1.65) {};
\draw[line width=1.15pt] (g5a)--(g5u);
\draw[line width=1.15pt] (g5a)--(g5v);
\draw[line width=1.15pt] (g5b)--(g5u);
\draw[line width=1.15pt] (g5b)--(g5v);
\draw[line width=1.15pt] (g5u)--(g5v);
\end{tikzpicture}\newpage
\noindent\begin{tikzpicture}[remember picture,overlay,x=1in,y=1in,shift={(current page.south west)}]
\node[anchor=north west,inner sep=0pt,align=left,text width=7.18in,font=\fontsize{13}{17}\selectfont] at (0.64,10.48) {Week 20 / Averaging / Grades 4--5};
\draw[line width=.5pt] (.64,10.13)--(7.86,10.13);
\node[anchor=north west,inner sep=0pt,align=left,text width=7.18in,font=\fontsize{9.5}{13.5}\selectfont] at (0.64,0.48) {Bellingham Math Circle / Week 20 / GA20-45-v3};
\node[anchor=north east,inner sep=0pt,font=\fontsize{9.5}{12}\selectfont] at (7.86,.48) {3};
\node[anchor=north west,inner sep=0pt,align=left,text width=7.18in,font=\fontsize{14.5}{18.5}\selectfont] at (0.64,9.86) {\textbf{Problem 3:} A board has finitely many shapes and lines. Every shape can reach a square along the lines. Explain why no circle number can be below the smallest square number it can reach or above the largest one. Your explanation must work for every such board.};
\node[draw,fill=white,rectangle,line width=1.05pt,minimum size=0.68in,inner sep=0pt,font=\fontsize{19}{25}\selectfont] (g6a) at (1.1,6.25) {0};
\node[draw,fill=white,circle,line width=1.05pt,minimum size=0.68in,inner sep=0pt,font=\fontsize{19}{25}\selectfont] (g6u) at (2.6,6.25) {};
\node[draw,fill=white,circle,line width=1.05pt,minimum size=0.68in,inner sep=0pt,font=\fontsize{19}{25}\selectfont] (g6v) at (4.15,6.25) {};
\node[draw,fill=white,circle,line width=1.05pt,minimum size=0.68in,inner sep=0pt,font=\fontsize{19}{25}\selectfont] (g6w) at (5.7,6.25) {};
\node[draw,fill=white,rectangle,line width=1.05pt,minimum size=0.68in,inner sep=0pt,font=\fontsize{19}{25}\selectfont] (g6b) at (7.25,6.25) {12};
\node[draw,fill=white,rectangle,line width=1.05pt,minimum size=0.68in,inner sep=0pt,font=\fontsize{19}{25}\selectfont] (g6c) at (4.15,4.65) {6};
\node[draw,fill=white,circle,line width=1.05pt,minimum size=0.68in,inner sep=0pt,font=\fontsize{19}{25}\selectfont] (g6t) at (4.15,7.8) {};
\draw[line width=1.15pt] (g6a)--(g6u);
\draw[line width=1.15pt] (g6u)--(g6v);
\draw[line width=1.15pt] (g6v)--(g6w);
\draw[line width=1.15pt] (g6w)--(g6b);
\draw[line width=1.15pt] (g6v)--(g6c);
\draw[line width=1.15pt] (g6u)--(g6t);
\draw[line width=1.15pt] (g6t)--(g6w);
\end{tikzpicture}\newpage
\noindent\begin{tikzpicture}[remember picture,overlay,x=1in,y=1in,shift={(current page.south west)}]
\node[anchor=north west,inner sep=0pt,align=left,text width=7.18in,font=\fontsize{13}{17}\selectfont] at (0.64,10.48) {Week 20 / Averaging / Grades 4--5};
\draw[line width=.5pt] (.64,10.13)--(7.86,10.13);
\node[anchor=north west,inner sep=0pt,align=left,text width=7.18in,font=\fontsize{9.5}{13.5}\selectfont] at (0.64,0.48) {Bellingham Math Circle / Week 20 / GA20-45-v3};
\node[anchor=north east,inner sep=0pt,font=\fontsize{9.5}{12}\selectfont] at (7.86,.48) {4};
\node[anchor=north west,inner sep=0pt,align=left,text width=7.18in,font=\fontsize{14.5}{18.5}\selectfont] at (0.64,9.86) {\textbf{Problem 4:} Fill the circles. Which circles can equal the largest square number on their own board? Explain what makes the two boards different.};
\node[draw,fill=white,rectangle,line width=1.05pt,minimum size=0.68in,inner sep=0pt,font=\fontsize{19}{25}\selectfont] (g7a) at (1.15,7.6) {0};
\node[draw,fill=white,circle,line width=1.05pt,minimum size=0.68in,inner sep=0pt,font=\fontsize{19}{25}\selectfont] (g7u) at (2.7,7.6) {};
\node[draw,fill=white,rectangle,line width=1.05pt,minimum size=0.68in,inner sep=0pt,font=\fontsize{19}{25}\selectfont] (g7b) at (4.25,7.6) {8};
\node[draw,fill=white,circle,line width=1.05pt,minimum size=0.68in,inner sep=0pt,font=\fontsize{19}{25}\selectfont] (g7v) at (5.8,7.6) {};
\node[draw,fill=white,circle,line width=1.05pt,minimum size=0.68in,inner sep=0pt,font=\fontsize{19}{25}\selectfont] (g7w) at (7.35,7.6) {};
\draw[line width=1.15pt] (g7a)--(g7u);
\draw[line width=1.15pt] (g7u)--(g7b);
\draw[line width=1.15pt] (g7b)--(g7v);
\draw[line width=1.15pt] (g7v)--(g7w);
\node[draw,fill=white,rectangle,line width=1.05pt,minimum size=0.68in,inner sep=0pt,font=\fontsize{19}{25}\selectfont] (g8a) at (1.3,4.3) {0};
\node[draw,fill=white,circle,line width=1.05pt,minimum size=0.68in,inner sep=0pt,font=\fontsize{19}{25}\selectfont] (g8u) at (3.25,4.3) {};
\node[draw,fill=white,circle,line width=1.05pt,minimum size=0.68in,inner sep=0pt,font=\fontsize{19}{25}\selectfont] (g8v) at (5.2,4.3) {};
\node[draw,fill=white,rectangle,line width=1.05pt,minimum size=0.68in,inner sep=0pt,font=\fontsize{19}{25}\selectfont] (g8b) at (7.15,4.3) {12};
\node[draw,fill=white,circle,line width=1.05pt,minimum size=0.68in,inner sep=0pt,font=\fontsize{19}{25}\selectfont] (g8w) at (5.2,2.55) {};
\draw[line width=1.15pt] (g8a)--(g8u);
\draw[line width=1.15pt] (g8u)--(g8v);
\draw[line width=1.15pt] (g8v)--(g8b);
\draw[line width=1.15pt] (g8v)--(g8w);
\end{tikzpicture}\newpage
\noindent\begin{tikzpicture}[remember picture,overlay,x=1in,y=1in,shift={(current page.south west)}]
\node[anchor=north west,inner sep=0pt,align=left,text width=7.18in,font=\fontsize{13}{17}\selectfont] at (0.64,10.48) {Week 20 / Averaging / Grades 4--5};
\draw[line width=.5pt] (.64,10.13)--(7.86,10.13);
\node[anchor=north west,inner sep=0pt,align=left,text width=7.18in,font=\fontsize{9.5}{13.5}\selectfont] at (0.64,0.48) {Bellingham Math Circle / Week 20 / GA20-45-v3};
\node[anchor=north east,inner sep=0pt,font=\fontsize{9.5}{12}\selectfont] at (7.86,.48) {5};
\node[anchor=north west,inner sep=0pt,align=left,text width=7.18in,font=\fontsize{14.5}{18.5}\selectfont] at (0.64,9.86) {\textbf{Problem 5:} Can the same board, with the same square numbers, have two different fillings that both follow the rule? Explain your answer for every finite board where each shape can reach a square.};
\node[draw,fill=white,rectangle,line width=1.05pt,minimum size=0.68in,inner sep=0pt,font=\fontsize{19}{25}\selectfont] (g9a) at (1.2,7.1) {0};
\node[draw,fill=white,circle,line width=1.05pt,minimum size=0.68in,inner sep=0pt,font=\fontsize{19}{25}\selectfont] (g9u) at (2.7,7.1) {};
\node[draw,fill=white,circle,line width=1.05pt,minimum size=0.68in,inner sep=0pt,font=\fontsize{19}{25}\selectfont] (g9v) at (4.4,7.85) {};
\node[draw,fill=white,circle,line width=1.05pt,minimum size=0.68in,inner sep=0pt,font=\fontsize{19}{25}\selectfont] (g9w) at (4.4,6.35) {};
\node[draw,fill=white,rectangle,line width=1.05pt,minimum size=0.68in,inner sep=0pt,font=\fontsize{19}{25}\selectfont] (g9b) at (6.15,7.85) {16};
\node[draw,fill=white,rectangle,line width=1.05pt,minimum size=0.68in,inner sep=0pt,font=\fontsize{19}{25}\selectfont] (g9c) at (6.15,6.35) {8};
\draw[line width=1.15pt] (g9a)--(g9u);
\draw[line width=1.15pt] (g9u)--(g9v);
\draw[line width=1.15pt] (g9u)--(g9w);
\draw[line width=1.15pt] (g9v)--(g9w);
\draw[line width=1.15pt] (g9v)--(g9b);
\draw[line width=1.15pt] (g9w)--(g9c);
\node[draw,fill=white,rectangle,line width=1.05pt,minimum size=0.68in,inner sep=0pt,font=\fontsize{19}{25}\selectfont] (g10a) at (1.2,3.5) {0};
\node[draw,fill=white,circle,line width=1.05pt,minimum size=0.68in,inner sep=0pt,font=\fontsize{19}{25}\selectfont] (g10u) at (2.7,3.5) {};
\node[draw,fill=white,circle,line width=1.05pt,minimum size=0.68in,inner sep=0pt,font=\fontsize{19}{25}\selectfont] (g10v) at (4.4,4.25) {};
\node[draw,fill=white,circle,line width=1.05pt,minimum size=0.68in,inner sep=0pt,font=\fontsize{19}{25}\selectfont] (g10w) at (4.4,2.75) {};
\node[draw,fill=white,rectangle,line width=1.05pt,minimum size=0.68in,inner sep=0pt,font=\fontsize{19}{25}\selectfont] (g10b) at (6.15,4.25) {16};
\node[draw,fill=white,rectangle,line width=1.05pt,minimum size=0.68in,inner sep=0pt,font=\fontsize{19}{25}\selectfont] (g10c) at (6.15,2.75) {8};
\draw[line width=1.15pt] (g10a)--(g10u);
\draw[line width=1.15pt] (g10u)--(g10v);
\draw[line width=1.15pt] (g10u)--(g10w);
\draw[line width=1.15pt] (g10v)--(g10w);
\draw[line width=1.15pt] (g10v)--(g10b);
\draw[line width=1.15pt] (g10w)--(g10c);
\end{tikzpicture}\newpage
\noindent\begin{tikzpicture}[remember picture,overlay,x=1in,y=1in,shift={(current page.south west)}]
\node[anchor=north west,inner sep=0pt,align=left,text width=7.18in,font=\fontsize{13}{17}\selectfont] at (0.64,10.48) {Week 20 / Averaging / Grades 4--5};
\draw[line width=.5pt] (.64,10.13)--(7.86,10.13);
\node[anchor=north west,inner sep=0pt,align=left,text width=7.18in,font=\fontsize{9.5}{13.5}\selectfont] at (0.64,0.48) {Bellingham Math Circle / Week 20 / GA20-45-v3};
\node[anchor=north east,inner sep=0pt,font=\fontsize{9.5}{12}\selectfont] at (7.86,.48) {6};
\node[anchor=north west,inner sep=0pt,align=left,text width=7.18in,font=\fontsize{14.5}{18.5}\selectfont] at (0.64,9.86) {\textbf{Problem 6:} Use whole numbers from 0 to 4. Find every filling of the upper board and of the lower board. How many fillings does each board have? Explain why none are missing.};
\node[draw,fill=white,circle,line width=1.05pt,minimum size=0.68in,inner sep=0pt,font=\fontsize{19}{25}\selectfont] (g11a) at (4.25,8.0) {};
\node[draw,fill=white,circle,line width=1.05pt,minimum size=0.68in,inner sep=0pt,font=\fontsize{19}{25}\selectfont] (g11b) at (2.8,6.05) {};
\node[draw,fill=white,circle,line width=1.05pt,minimum size=0.68in,inner sep=0pt,font=\fontsize{19}{25}\selectfont] (g11c) at (5.7,6.05) {};
\draw[line width=1.15pt] (g11a)--(g11b);
\draw[line width=1.15pt] (g11b)--(g11c);
\draw[line width=1.15pt] (g11c)--(g11a);
\draw[rounded corners=.18in,line width=.7pt,dashed] (.68,2.65) rectangle (7.82,4.35);
\node[draw,fill=white,circle,line width=1.05pt,minimum size=0.68in,inner sep=0pt,font=\fontsize{19}{25}\selectfont] (g12a) at (1.2,3.5) {};
\node[draw,fill=white,circle,line width=1.05pt,minimum size=0.68in,inner sep=0pt,font=\fontsize{19}{25}\selectfont] (g12b) at (3.0,3.5) {};
\node[draw,fill=white,circle,line width=1.05pt,minimum size=0.68in,inner sep=0pt,font=\fontsize{19}{25}\selectfont] (g12c) at (5.5,3.5) {};
\node[draw,fill=white,circle,line width=1.05pt,minimum size=0.68in,inner sep=0pt,font=\fontsize{19}{25}\selectfont] (g12d) at (7.3,3.5) {};
\draw[line width=1.15pt] (g12a)--(g12b);
\draw[line width=1.15pt] (g12c)--(g12d);
\end{tikzpicture}\newpage
\noindent\begin{tikzpicture}[remember picture,overlay,x=1in,y=1in,shift={(current page.south west)}]
\node[anchor=north west,inner sep=0pt,align=left,text width=7.18in,font=\fontsize{13}{17}\selectfont] at (0.64,10.48) {Week 20 / Averaging / Grades 4--5};
\draw[line width=.5pt] (.64,10.13)--(7.86,10.13);
\node[anchor=north west,inner sep=0pt,align=left,text width=7.18in,font=\fontsize{9.5}{13.5}\selectfont] at (0.64,0.48) {Bellingham Math Circle / Week 20 / GA20-45-v3};
\node[anchor=north east,inner sep=0pt,font=\fontsize{9.5}{12}\selectfont] at (7.86,.48) {7};
\node[anchor=north west,inner sep=0pt,align=left,text width=7.18in,font=\fontsize{14.5}{18.5}\selectfont] at (0.64,9.86) {\textbf{Problem 7:} Which boards can you fill with whole numbers? Fill the other boards using the fraction cards.};
\node[draw,fill=white,rectangle,line width=1.05pt,minimum size=0.68in,inner sep=0pt,font=\fontsize{19}{25}\selectfont] (g130) at (1.4,8.15) {0};
\node[draw,fill=white,circle,line width=1.05pt,minimum size=0.68in,inner sep=0pt,font=\fontsize{19}{25}\selectfont] (g131) at (4.25,8.15) {};
\node[draw,fill=white,rectangle,line width=1.05pt,minimum size=0.68in,inner sep=0pt,font=\fontsize{19}{25}\selectfont] (g132) at (7.1,8.15) {1};
\draw[line width=1.15pt] (g130)--(g131);
\draw[line width=1.15pt] (g131)--(g132);
\node[draw,fill=white,rectangle,line width=1.05pt,minimum size=0.68in,inner sep=0pt,font=\fontsize{19}{25}\selectfont] (g140) at (1.4,6.65) {0};
\node[draw,fill=white,circle,line width=1.05pt,minimum size=0.68in,inner sep=0pt,font=\fontsize{19}{25}\selectfont] (g141) at (3.3,6.65) {};
\node[draw,fill=white,circle,line width=1.05pt,minimum size=0.68in,inner sep=0pt,font=\fontsize{19}{25}\selectfont] (g142) at (5.2,6.65) {};
\node[draw,fill=white,rectangle,line width=1.05pt,minimum size=0.68in,inner sep=0pt,font=\fontsize{19}{25}\selectfont] (g143) at (7.1,6.65) {3};
\draw[line width=1.15pt] (g140)--(g141);
\draw[line width=1.15pt] (g141)--(g142);
\draw[line width=1.15pt] (g142)--(g143);
\node[draw,fill=white,rectangle,line width=1.05pt,minimum size=0.68in,inner sep=0pt,font=\fontsize{19}{25}\selectfont] (g150) at (1.4,5.15) {0};
\node[draw,fill=white,circle,line width=1.05pt,minimum size=0.68in,inner sep=0pt,font=\fontsize{19}{25}\selectfont] (g151) at (4.25,5.15) {};
\node[draw,fill=white,rectangle,line width=1.05pt,minimum size=0.68in,inner sep=0pt,font=\fontsize{19}{25}\selectfont] (g152) at (7.1,5.15) {2};
\draw[line width=1.15pt] (g150)--(g151);
\draw[line width=1.15pt] (g151)--(g152);
\node[draw,fill=white,rectangle,line width=1.05pt,minimum size=0.68in,inner sep=0pt,font=\fontsize{19}{25}\selectfont] (g160) at (1.4,3.65) {0};
\node[draw,fill=white,circle,line width=1.05pt,minimum size=0.68in,inner sep=0pt,font=\fontsize{19}{25}\selectfont] (g161) at (3.3,3.65) {};
\node[draw,fill=white,circle,line width=1.05pt,minimum size=0.68in,inner sep=0pt,font=\fontsize{19}{25}\selectfont] (g162) at (5.2,3.65) {};
\node[draw,fill=white,rectangle,line width=1.05pt,minimum size=0.68in,inner sep=0pt,font=\fontsize{19}{25}\selectfont] (g163) at (7.1,3.65) {1};
\draw[line width=1.15pt] (g160)--(g161);
\draw[line width=1.15pt] (g161)--(g162);
\draw[line width=1.15pt] (g162)--(g163);
\draw[line width=.65pt] (1.3699999999999999,1.15) rectangle (2.9299999999999997,2.3);
\draw[fill=black!30] (1.5499999999999998,1.82) rectangle (2.15,2.08);
\draw[fill=white] (2.15,1.82) rectangle (2.75,2.08);
\node[font=\fontsize{18}{22}\selectfont] at (2.15,1.48) {$\frac{1}{2}$};
\draw[line width=.65pt] (3.4699999999999998,1.15) rectangle (5.03,2.3);
\draw[fill=black!30] (3.65,1.82) rectangle (4.05,2.08);
\draw[fill=white] (4.05,1.82) rectangle (4.45,2.08);
\draw[fill=white] (4.45,1.82) rectangle (4.85,2.08);
\node[font=\fontsize{18}{22}\selectfont] at (4.25,1.48) {$\frac{1}{3}$};
\draw[line width=.65pt] (5.569999999999999,1.15) rectangle (7.13,2.3);
\draw[fill=black!30] (5.75,1.82) rectangle (6.15,2.08);
\draw[fill=black!30] (6.15,1.82) rectangle (6.55,2.08);
\draw[fill=white] (6.55,1.82) rectangle (6.95,2.08);
\node[font=\fontsize{18}{22}\selectfont] at (6.35,1.48) {$\frac{2}{3}$};
\end{tikzpicture}\null\end{document}